% Answers to selected exercises, chapter 11 (high dimensions and hubs).
% Every number is recomputed by checks/ch11_examples.py.

\paragraph{Exercise \ref{exr:ch11:scale}}
Write $x=\sigma g$ with $g\sim\Normal(0,I_d)$. Then $\norm x^2=\sigma^2\norm g^2$, so
$\E\norm x^2=d\sigma^2$ and $\Var\norm x^2=\sigma^4\cdot2d$. The relative spread is
$\sqrt{2d}\,\sigma^2/(d\sigma^2)=\sqrt{2/d}$, the same as for $\sigma=1$. Concentration
is a statement about ratios, and a change of units cannot affect it.

\paragraph{Exercise \ref{exr:ch11:sum}}
Count the pairs $(q,x)$ with $x\in\mathrm{top}_k(q)$ in two ways. Grouping by query,
each $q$ contributes exactly $k$ pairs, for $k|Q|$ in total. Grouping by stored point,
$x$ contributes $N_k(x)$ pairs. So $\sum_xN_k(x)=k|Q|$, and in the corpus-to-corpus
experiment the mean count is $k$ for every corpus, metric and dimension. The mean
therefore carries no information about hubness, and a statistic of shape, such as the
skewness, the largest count or the Robin Hood index, is needed.

\paragraph{Exercise \ref{exr:ch11:shift}}
$\norm{(x+b)-(q+b)}=\norm{x-q}$, so every distance, every neighbour list and every count
is unchanged. The inner product is not shift-invariant,
$\inner{q+b}{x+b}=\inner qx+\inner bx+\inner qb+\norm b^2$, and the term $\inner bx$
favours the same stored points for every query. Cosine similarity changes as well.
That is why centering is a remedy for inner-product and cosine hubs and does nothing
for Euclidean ones.

\paragraph{Exercise \ref{exr:ch11:nullskew}}
Under the no-structure null of \cref{eq:ch11:null} the counts are close to
$\mathrm{Poisson}(10)$, whose skewness is $1/\sqrt{10}=0.316$. A reading of $0.30$ is
what random assignment of slots produces. It gives no evidence of hubness, mild or
otherwise, and the corpus may even be slightly more even than chance. It does not
prove that hubs are absent, since one skewness value summarizes the whole count
distribution. The largest count against the ceiling $c^\star$ of
\cref{eq:ch11:ceiling} is the more direct check.

\paragraph{Exercise \ref{exr:ch11:shell}}
$1-0.98^{50}=0.6358$. The lower bound is $1-e^{-0.02\cdot50}=1-e^{-1}=0.6321$, close
to the exact value because $\epsilon$ is small.

\paragraph{Exercise \ref{exr:ch11:five2}}
The two nearest neighbours are $A\colon B,C$, $B\colon A,D$, $C\colon A,D$,
$D\colon A,B$ and $E\colon A,B$, using $BD=3.081<BE=3.499<BC=3.8$ and
$CD=3.202<CE=3.606$. So $N_2=(4,3,1,2,0)$ for $(A,B,C,D,E)$, with sum $10$ and mean $2$.
The deviations are $2,1,-1,0,-2$, the variance is $10/5=2$, and the cubed deviations
sum to $8+1-1+0-8=0$, so the skewness is $0$. The Robin Hood index is
$(2+1+1+0+2)/(2\cdot10)=0.3$, against $0.6$ for $N_1$. The point $A$ is still the
busiest, but at $k=2$ the five points share the slots symmetrically around the mean.

\paragraph{Exercise \ref{exr:ch11:jlk}}
$\epsilon^2/2-\epsilon^3/3=0.005-0.000333=0.004667$ and $4\ln10^4=36.84$, so
$m\ge7894.6$ and $m=7895$ suffices. The tolerance enters as $1/\epsilon^2$ and the
number of points only as $\ln n$.

\paragraph{Exercise \ref{exr:ch11:ceil}}
The Poisson skewness is $1/\sqrt{20}=0.224$. With $\mu=20$ and $n=5000$,
$5000\Pr[\mathrm{Poisson}(20)\ge38]=1.09$ and $5000\Pr[\mathrm{Poisson}(20)\ge39]=0.54$,
so $c^\star=38$. A busiest count of $190$ is a hub excess of $190/38=5$.

\paragraph{Exercise \ref{exr:ch11:model}}
$c=\Phi^{-1}(0.01)=-2.326$. At $z=0$, $p=\Phi(\sqrt6\,c/2)=\Phi(-2.849)=0.00219$, an
expected count of $2.19$. At $z=-1$, $p=\Phi\big((\sqrt6\,c+\sqrt2)/2\big)=\Phi(-2.142)=0.0161$,
an expected count of $16.1$. One standard deviation of centrality multiplies the
expected count by more than seven.

\paragraph{Exercise \ref{exr:ch11:prog-hub}}
For the embedding part, $\norm{Ux-Uy}=\norm{x-y}$ because $U\T U=I_5$, so every
distance, every neighbour list and every count is the same in $\R^{500}$ as in
$\R^5$. The check script confirms the counts are identical. Hubness depends on the
geometry of the data, not on the number of coordinates used to write them down, so
for data on or near a low-dimensional set it is the intrinsic dimension that matters.
